\section*{Chapitre \ref{ch:b2:coordination-complexes} --- Métaux de transition et complexes}

\begin{solution}{exo:b2:coordination-complexes:1}
\ce{Ti^3+} $d^1$, \ce{V^3+} $d^2$, \ce{Cr^3+} $d^3$, \ce{Mn^2+} $d^5$, \ce{Fe^3+} $d^5$, \ce{Co^2+}
$d^7$, \ce{Ni^2+} $d^8$, \ce{Zn^2+} $d^{10}$.
\end{solution}

\begin{solution}{exo:b2:coordination-complexes:2}
Eau 1, en 2, éthanedioate 2, edta 6, thiocyanate 1~; le thiocyanate est ambidenté
(par S ou par N).
\end{solution}

\begin{solution}{exo:b2:coordination-complexes:3}
Sulfate de tétraamminecuivre(II)~; hexacyanidoferrate(III) de potassium~;
chlorure de tétraamminedichloridocobalt(III)~; tétracarbonylnickel(0).
\end{solution}

\begin{solution}{exo:b2:coordination-complexes:4}
Linéaire~; tétraédrique~; plan carré ($d^8$, troisième série)~; octaédrique.
\end{solution}

\begin{solution}{exo:b2:coordination-complexes:5}
Deux~: cis (les deux chlorures adjacents) et trans (opposés). Dans un tétraèdre, il n'y en
aurait qu'un.
\end{solution}

\begin{solution}{exo:b2:coordination-complexes:6}
\ce{Cr(CO)6} 18~; \ce{Fe(CO)5} 18~; \ce{Ni(CO)4} 18~; \ce{V(CO)6}~: $5 + 12 = 17$~;
\ce{Mn2(CO)10} 18 par manganèse. \ce{V(CO)6} déroge à la règle~: il est \omterm{def:b2:diatomic-mos:magnetism}{paramagnétique} et
facilement réduit en \ce{[V(CO)6]-}, qui en a 18.
\end{solution}

\begin{solution}{exo:b2:coordination-complexes:7}
Ferrocène~: Fe(II) $d^6$ + 12 = 18. Cobaltocène~: Co(II) $d^7$ + 12 = 19. Le cobaltocène perd
facilement son électron excédentaire, ce qui donne l'ion cobaltocénium \ce{[Co(C5H5)2]+}, à 18 électrons.
\end{solution}

\begin{solution}{exo:b2:coordination-complexes:8}
Un edta remplace six molécules d'eau~: \ce{[Ni(H2O)6]^2+ + edta^4- -> [Ni(edta)]^2- +
6H2O}, sept particules à partir de deux. L'\omterm{def:b2:reaction-free-energy:reaction-entropy}{entropie de réaction} est fortement positive, et
$K$ très grande. Six ligands séparés remplacent six molécules d'eau sans gain sur le
nombre de particules.
\end{solution}

\begin{solution}{exo:b2:coordination-complexes:9}
Nitrito-$\kappa N$ (\ce{Co-NO2}, l'isomère nitro, jaune) et nitrito-$\kappa O$
(\ce{Co-O-N=O}, l'isomère nitrito, rouge).
\end{solution}

\begin{solution}{exo:b2:coordination-complexes:10}
Trois~: l'isomère trans (achiral~: plans de symétrie) et les deux énantiomères de l'isomère
cis (les deux chélates et les deux chlorures s'enroulent dans un sens ou dans l'autre).
\end{solution}

\begin{solution}{exo:b2:coordination-complexes:11}
Hexagone plan~: B en positions 1,2 (ortho), 1,3 (méta), 1,4 (para)~: trois isomères.
Prisme trigonal~: B le long d'une arête d'un triangle, le long d'une arête verticale, ou en diagonale d'une
face rectangulaire~: trois. L'octaèdre en donne deux. N'en trouver jamais plus de deux, malgré
de nombreuses tentatives et plusieurs composés, exclut les deux autres géométries --- à condition qu'aucun
troisième isomère n'ait échappé à l'observation.
\end{solution}

\begin{solution}{exo:b2:coordination-complexes:12}
\ce{[RhCl(PPh3)3]}~: Rh(I) $d^8$ + 4 × 2 = 16. \ce{[IrH(CO)(PPh3)3]}~: Ir(I) $d^8$ + 5 × 2 =
18. Le complexe du rhodium, à 16 électrons, peut en accepter deux de plus (addition oxydante,
\cref{ch:b2:catalytic-cycles}).
\end{solution}

\begin{solution}{pb:b2:coordination-complexes:1}
\textbf{1.} Les chlorures extérieurs à la \omterm{def:b2:coordination-complexes:coordination-sphere}{sphère de coordination}, libres en solution.
\textbf{2.} A~: 0~; B~: 1~; C et D~: 2.
\textbf{3.} A~: $[\ce{Co(NH3)6}]^{3+}$ + 3 \ce{Cl-}~: 4 ions~; B~: 3~; C, D~: 2. Les valeurs concordent.
\textbf{4.} +3.
\textbf{5.} $d^6$.
\textbf{6.} Six dans chacun~: 6 \ce{NH3}~; 5 \ce{NH3} + 1 Cl~; 4 \ce{NH3} + 2 Cl.
\textbf{7.} \ce{[Co(NH3)6]Cl3}~; \ce{[CoCl(NH3)5]Cl2}~; \ce{[CoCl2(NH3)4]Cl} (deux fois).
\textbf{8.} Chlorure d'hexaamminecobalt(III)~; chlorure de pentaamminechloridocobalt(III).
\textbf{9.} Tétraamminedichloridocobalt(III).
\textbf{10.} Des isomères géométriques (cis--trans).
\textbf{11.} \ce{[CoCl3(NH3)3]}, un complexe neutre~: pas d'ions, aucun chlorure précipité
d'emblée.
\textbf{12.} Deux~: fac et mer.
\textbf{13.} Cis~: les deux Cl à $90^\circ$~; trans~: à $180^\circ$.
\textbf{14.} Trois (positions ortho, méta et para sur l'hexagone).
\textbf{15.} Trois.
\textbf{16.} Cela appuie l'octaèdre, la seule des trois géométries qui en prévoit deux.
\textbf{17.} Non~: l'absence d'un troisième isomère est une preuve négative. Un troisième isomère aurait
exclu l'octaèdre.
\textbf{18.} Celui qui forme le chélate~: un éthanedioate bidenté ne peut occuper que deux
positions cis.
\textbf{19.} Trois chélates en autour du cobalt, chacun occupant deux positions cis.
\textbf{20.} Ni l'un ni l'autre~: pas de plan de symétrie et pas de centre d'inversion.
\textbf{21.} Deux énantiomères, $\Delta$ et $\Lambda$.
\textbf{22.} Que les six ligands sont disposés dans les trois dimensions, comme dans un
octaèdre~; une disposition plane donnerait un ion achiral.
\textbf{23.} Cis~: chiral (deux énantiomères). Trans~: achiral.
\textbf{24.} Il montra que l'activité optique tient à la géométrie moléculaire en général,
et non au carbone, et confirma l'octaèdre.
\textbf{25.} L'octaèdre prévoit $\boldsymbol{2}$ isomères de \ce{[CoCl2(NH3)4]+}~; l'hexagone
plan et le prisme trigonal en prévoient $\boldsymbol{3}$.
\end{solution}
